\chapter{Introduction}

\section{Introduction}
Every organization produces and receives legal documents: non-disclosure agreements (NDAs) with
partners, service and supply contracts, office leases, employment and appointment letters,
memoranda of understanding, affidavits and powers of attorney. These documents decide who owes
what to whom, when obligations begin and end, and how much notice must be given to leave an
agreement. Alongside them sit \emph{reference} documents such as the Contract Act 1872, company
policies and handbooks, which explain the rules that apply to those agreements.

In most small and medium organizations in Pakistan these documents live in shared folders, e-mail
attachments and filing cabinets. Many are scanned, and some are handwritten on stamp paper. When a
manager needs a simple fact, for example the date an NDA was signed or the notice period of a lease,
someone has to find the right file and read it. Keyword search inside a PDF reader cannot help with
scanned pages, cannot understand a question phrased in different words from the document, and
cannot answer questions that span many documents, such as \emph{``Which contracts expire in 2027?''}.

Recent AI document assistants allow questions in natural language. They are usually built on
Retrieval-Augmented Generation (RAG)~\cite{lewis2020rag}, in which passages relevant to a question
are retrieved and passed to a Large Language Model (LLM) so that the answer is grounded in the text.
However, general-purpose tools treat every upload as loose text. They do not keep a structured
register of records, do not know the difference between a contract and a law book, do not control
which employee may see which document, keep no history of who changed what, and cannot prove where
a date or party name came from.

The first version of this project was such a general ``chat with your PDF'' tool for individual
users. Following the evaluation of that version, the project was redesigned around a single domain,
the legal records of an organization, to address the gaps above. This report presents the
redesigned system, \textbf{\sysname{}}: a web-based legal records intelligence system in which the
employees of an organization share one access-controlled document repository, every record is
turned into verifiable structured facts, questions are answered with citations that open the source
page, and every action is kept in history. \sysname{} was developed as a Final Year Project at
COMSATS University Islamabad, Vehari Campus, under the supervision of Dr.\ Umar Rashid.

\section{Vision Statement}
\sysname{} is intended for organizations that hold legal documents and need to find facts in them
quickly and safely: administration, human resources, procurement and legal staff, and the managers
and auditors who supervise them.

The vision of \sysname{} is to turn an organization's pile of contracts, scans and law books into a
\emph{trustworthy, queryable records system}: one in which any employee can ask a question in plain
English and receive an answer that is backed by the exact page it came from, limited to the
documents that employee is allowed to see, and recorded in a history that can be audited later.

The innovation of the system lies in combining, in one platform, (i) reading of digital, scanned and
handwritten pages, (ii) automatic separation of records from reference material, (iii) extraction of
structured record fields in which every value is proved by a verbatim quote and checked against the
page, (iv) a query router that answers fact questions from a database table and explanation
questions from retrieved passages, and (v) access control applied \emph{before} retrieval, so the
language model never sees a document the user may not open.

\section{Related System Analysis}
Two groups of existing systems were analyzed: general AI document assistants and commercial contract
analysis platforms. The descriptions below are based on publicly available information about these
products.

\textbf{ChatPDF}~\cite{chatpdf} and \textbf{Google NotebookLM}~\cite{notebooklm} let a user upload
files and ask questions about them with cited answers. They need almost no setup, but they are
built for one person's study material: there is no shared organizational repository with roles, no
structured register of contract dates and parties, no document versioning and no audit trail.
Documents are processed on the vendor's servers under the vendor's policies.

\textbf{Kira}~\cite{kira} and \textbf{Luminance}~\cite{luminance} are commercial machine-learning
platforms used by law firms to review contracts, find clauses and extract provisions. They are
powerful, but they are priced for large legal teams, are not designed for handwritten Pakistani stamp
paper documents, and are not available to small organizations at an academic or SME budget.

\textbf{Contract lifecycle management (CLM) systems} such as Ironclad~\cite{ironclad} manage the
drafting, approval and renewal of contracts created \emph{inside} the system. They store metadata
well, but they assume structured, digital contracts and offer limited help with legacy scans or with
questions answered from a law book.

Table~\ref{tab:related} summarizes the weaknesses identified and how \sysname{} addresses them.

\begin{table}[H]\centering
\caption{Comparative Analysis of Existing Systems}\label{tab:related}
\begin{tabularx}{\textwidth}{|p{3cm}|X|X|}\hline
\textbf{Application} & \textbf{Identified Weaknesses} & \textbf{Proposed Improvements} \\ \hline
ChatPDF / NotebookLM &
\begin{bullets}
\item Built for one user's study files
\item No roles, departments or sharing rules
\item No structured record of dates and parties
\item No version history or audit log
\end{bullets} &
\begin{bullets}
\item Shared organizational repository with four roles and department / named-person access
\item Structured records register with evidence for every value
\item Four kinds of history kept
\end{bullets} \\ \hline
Kira / Luminance &
\begin{bullets}
\item Enterprise pricing for large legal teams
\item Not designed for handwritten stamp-paper documents
\end{bullets} &
\begin{bullets}
\item Built on open-source components; free-tier LLM APIs or a local model
\item Handwriting recognition with a vision-language model
\end{bullets} \\ \hline
CLM systems (e.g.\ Ironclad) &
\begin{bullets}
\item Assume digital contracts created in the system
\item Little support for legacy scans or law-book questions
\end{bullets} &
\begin{bullets}
\item Ingests existing PDFs, DOCX files, scans and photos
\item Answers questions from reference material with page citations
\end{bullets} \\ \hline
\end{tabularx}
\end{table}

\section{Project Deliverables}
The tangible outputs of the project are listed below.
\begin{enumerate}
\item \textbf{\sysname{} Web Application (Frontend)} -- A Vue.js 3 single-page application with the
Landing page, Sign-up and Login, the Assistant (question answering), Documents, Document detail with
an embedded PDF viewer, Records register, Review queue, Team management, Audit log and Profile
screens.
\item \textbf{Backend Services} -- A Django 6 / Django REST Framework backend exposing REST APIs for
authentication, organizations and members, documents and versions, records and reviews, the
assistant, notifications and the audit log.
\item \textbf{Document Processing Pipeline} -- A database-backed job queue and worker that extract
text (digital, OCR and handwriting), classify the document, chunk it by its section structure,
compute embeddings, label clauses and extract record fields with evidence.
\item \textbf{Legal Assistant} -- A query router and three answer engines (lookup, knowledge and
search) built on hybrid retrieval with Reciprocal Rank Fusion, with citations that open the source
page.
\item \textbf{Fine-tuned Clause Classifier} -- A transformer model fine-tuned on LEDGAR to label
contract clauses with 100 clause types, with training scripts, baselines and a Google Colab notebook.
\item \textbf{Evaluation Suite} -- Reproducible scripts that measure retrieval, field extraction,
answer quality, query routing, handwriting recognition and clause classification, and write the
results into one report.
\item \textbf{Database} -- A PostgreSQL 16 database with the pgvector extension, an HNSW vector index
and a GIN full-text index.
\item \textbf{Deployment Configuration} -- Dockerfiles and a Docker Compose file for the database,
backend, worker and web server.
\item \textbf{Project Documentation} -- The SRS and SDD (Chapters 3 and 4), the test and evaluation
results (Chapter 6) and this report.
\end{enumerate}

\section{System Limitations / Constraints}
\begin{itemize}
\item \textbf{Language} -- Documents and handwriting are supported in English only. Urdu documents are
outside the scope.
\item \textbf{File types and size} -- PDF, DOCX, PNG and JPG/JPEG files up to 25\,MB each are accepted.
\item \textbf{Handwriting quality} -- Recognition of handwriting depends on the legibility and resolution
of the scan; uncertain values are sent to the review queue instead of being trusted.
\item \textbf{External AI services} -- Classification, field extraction and written answers use a hosted
LLM (Groq, with Google Gemini as a fallback). Free tiers limit tokens per minute and per day, so heavy
use can slow answers. Search, lookup answers and access control keep working without the LLM.
\item \textbf{Not legal advice} -- The assistant reports what documents say; it does not give legal
advice, and its extracted values should be verified before important decisions.
\item \textbf{One organization per account} -- Each user belongs to exactly one organization.
\item \textbf{Platform} -- \sysname{} is a web application; no native mobile application is provided.
\item \textbf{Resources} -- As an academic project, large-scale deployment and load testing with many
concurrent organizations were outside the scope; the clause model was trained on a free Google Colab GPU.
\end{itemize}

\section{Tools and Technologies}
The technology stack separates the user interface, the application logic and the data, and keeps
all document processing on the server. Vue.js~\cite{vue} was chosen for the frontend because its
reactive components suit an interactive document viewer and assistant. Django~\cite{django} with
Django REST Framework~\cite{drf} provides authentication, an ORM, migrations and a mature Python
ecosystem for document processing and machine learning. PostgreSQL~\cite{postgres} stores all
relational data, and its pgvector extension~\cite{pgvector} stores embeddings in the same database
with HNSW indexing~\cite{malkov2018hnsw}, so that permission filters and vector search run in a
single query. Tesseract~\cite{smith2007tesseract} reads printed scans; a vision-language model
reads handwriting. Sentence embeddings are computed locally with
all-MiniLM-L6-v2~\cite{reimers2019sbert,wang2020minilm}, and the clause classifier is trained with
Hugging Face Transformers~\cite{wolf2020transformers}.

\begin{table}[H]\centering
\caption{Tools and Technologies Used in the Project}\label{tab:tools}
\begin{tabularx}{\textwidth}{|p{4.2cm}|p{2.8cm}|X|}\hline
\textbf{Tool / Technology} & \textbf{Layer} & \textbf{Purpose} \\ \hline
Vue.js 3, Vue Router, Vite & Frontend & Single-page application and build tooling \\ \hline
Tailwind CSS 4, Reka UI, Lucide & Frontend & Styling, accessible UI primitives and icons \\ \hline
vue-pdf-embed & Frontend & In-browser PDF viewer with evidence highlighting \\ \hline
Django 6, Django REST Framework & Backend & REST APIs, ORM, permissions, throttling \\ \hline
SimpleJWT & Backend & Token-based authentication \\ \hline
PostgreSQL 16 + pgvector & Database & Relational data, vector search (HNSW), full-text search (GIN) \\ \hline
pypdfium2, python-docx & Backend & PDF rendering and text extraction, DOCX reading \\ \hline
Tesseract OCR & Backend & Printed text recognition from scans \\ \hline
sentence-transformers (all-MiniLM-L6-v2) & Backend (local) & 384-dimensional sentence embeddings \\ \hline
Hugging Face Transformers, PyTorch & Machine learning & Fine-tuning the clause classifier \\ \hline
Groq API (gpt-oss-20b), Google Gemini & External service & LLM for routing, extraction and answers; vision model for handwriting \\ \hline
Docker, Docker Compose, Nginx, Gunicorn & Deployment & Containerized deployment \\ \hline
Git, GitHub & Development & Version control \\ \hline
\end{tabularx}
\end{table}

\section{Relevance to Course Modules}
The project integrates knowledge from several courses of the degree program.

\subsection{Programming Fundamentals}
The system is organized into Django applications with clear responsibilities (\code{accounts},
\code{documents}, \code{assistant}, \code{audit}, \code{notifications}) and an \code{ai} package that
wraps every model call. Domain concepts such as Document, DocumentVersion, Chunk, RecordField and
AuditEvent are modeled as classes, and the processing pipeline is split into small, separately
testable steps.

\subsection{Data Structures and Algorithms}
Retrieval uses an HNSW graph index for approximate nearest-neighbor search and an inverted index for
full-text search, combined with Reciprocal Rank Fusion~\cite{cormack2009rrf}. The chunker keeps a stack
of open section headings to build each chunk's section path. The job queue uses row-level locking
(\code{SELECT \dots{} FOR UPDATE SKIP LOCKED}) so several workers never take the same job.

\subsection{Database Management Systems}
The schema is normalized to Third Normal Form, with versioned documents, page and chunk tables,
uniqueness constraints (one value per field per version, one notification per reminder key), indexes
chosen for the main queries, and database-level append-only behavior for the audit log.

\subsection{Software Engineering}
The project followed an iterative Agile process. Requirements were elicited from the supervisor and
documented with use cases, functional and non-functional requirements and a traceability matrix. The
design uses UML activity, class, sequence and state diagrams, data flow diagrams and an ERD. Testing
combines automated unit and integration tests, black-box system tests and a quantitative evaluation.

\subsection{Other Relevant Courses}
Artificial intelligence and machine learning concepts (embeddings, transformers, fine-tuning,
evaluation metrics such as precision, recall, F1 and MRR) form the core of the system. Information
security concepts (role-based access control~\cite{sandhu1996rbac}, token authentication, rate limiting
and audit trails) protect the documents. Human-Computer Interaction principles guided the interface,
in particular showing evidence next to every extracted value.

\section{Chapter Summary}
This chapter introduced the problem of managing an organization's legal documents, the vision of
\sysname{}, and a comparison with existing systems. Project deliverables, constraints, the selected
technologies and the academic relevance of the project were also discussed. The next chapter defines
the problem precisely and describes the proposed solution and its modules.
